\documentclass[conference]{IEEEtran}

% ── Packages ────────────────────────────────────────────────────────────────
\usepackage[utf8]{inputenc}
\usepackage[T1]{fontenc}
\usepackage[english]{babel}
\usepackage{graphicx}
\usepackage{hyperref}
\usepackage{booktabs}
\usepackage{array}
\usepackage{amsmath}
\usepackage{amssymb}
\usepackage{cite}
\usepackage{enumitem}
\usepackage{xcolor}

\hypersetup{colorlinks=true, linkcolor=blue, urlcolor=blue, citecolor=blue}

% ════════════════════════════════════════════════════════════════════════════
\begin{document}

\title{A Custom Autonomy Stack for a Warehouse AMR%
\thanks{Source code:
\href{https://github.com/JordanPalafox/SmallAutonomousMobileRobot}{%
\texttt{github.com/JordanPalafox/SmallAutonomousMobileRobot}}}}

\author{
\IEEEauthorblockN{
Juan José Jáuregui Barba,
Victor Alejandro Meneses Garza,
Rosendo De Los Rios Moreno,\\
Jordan Palafox Salinas, and
Hugo Daniel Castillo Ovando
}
\IEEEauthorblockA{
\textit{Tecnológico de Monterrey} --- TE3003B, FJ2026\\
\{A00836722, A01384002, A01198515,
A00835705, A00836025\}@tec.mx
}
}
\maketitle

% ── Abstract ─────────────────────────────────────────────────────────────────
\begin{abstract}
This report documents the complete final system of the TE3003B
\textit{Puzzlebot Challenge} project: an autonomous forklift AMR that
executes full pick-and-place warehouse missions on a \mbox{3.65 × 4.85 m}
scaled track. The stack is entirely custom-built --- no Nav2, no Whisper,
no \texttt{robot\_localization}. The core subsystems are:
(1)~a C++ MCL/AMCL SLAM node with optional CUDA scoring, ICP
scan-to-scan and scan-to-map refinement, and a graph-SLAM back-end for
loop closure;
(2)~an ArUco beacon localization layer (\texttt{aruco\_localization}) that
rescues the particle filter when it drifts in low-feature corridors;
(3)~A*~+~Bug1 navigation with pure-pursuit trajectory tracking;
(4)~a YASMIN state machine that orchestrates full pick-and-place missions
across three zone types (rollers, racks, trucks);
(5)~a YOLOv8 logo classifier and QR-based docking for visual alignment at
pick/place sites;
(6)~an MFCC~+~VQ~+~HMM voice pipeline for command dispatch;
(7)~a Flask~+~SocketIO web dashboard with live telemetry and camera
streaming;
(8)~a Tang~Nano~20K FPGA lifter controlled via SPI from the Jetson Nano;
and (9)~a live Gazebo digital twin that mirrors the physical robot in
real time.
Two full pick-and-place missions were validated end-to-end on hardware:
Mission~1 (Rollers → Truck) and Mission~2 (Racks → Truck).
\end{abstract}

\begin{IEEEkeywords}
ROS~2, Puzzlebot, MCL, AMCL, CUDA, ICP, graph SLAM, loop closure,
ArUco, localization, A*, Bug1, pure-pursuit, YASMIN, FPGA, SPI,
YOLOv8, logo classifier, HMM, voice, Flask, SocketIO, digital twin,
Gazebo
\end{IEEEkeywords}

% ════════════════════════════════════════════════════════════════════════════
\section{Introduction}

The \textit{Puzzlebot Challenge} (TE3003B, ITESM FJ2026) asks teams to
build a fully autonomous forklift AMR that picks pallets from one zone
and delivers them to another inside a scaled warehouse. The arena
contains three zone types:
\textbf{Trucks} (loading docks, 3 units),
\textbf{Racks} (two-level shelving), and
\textbf{Rollers} (roller-table staging areas).

The team's principal constraint, maintained from Sprint~1 to the final
demonstration, was to implement every component from scratch:
\textbf{Nav2 is not used, Whisper is not used, \texttt{robot\_localization}
is not used}. The system is distributed across two compute nodes:
a \textbf{Jetson Nano 4~GB} on-board (sensing, SLAM, lifting) and a
\textbf{laptop} off-board (navigation, mission control, dashboard,
voice, digital twin), connected over Wi-Fi through shared DDS with
\texttt{ROS\_DOMAIN\_ID}.

Figure~\ref{fig:pista} shows the physical competition track and the
Puzzlebot.

\begin{figure}[htbp]
  \centering
  \includegraphics[width=\columnwidth]{assets/pista.jpeg}
  \caption{The \mbox{3.65 × 4.85 m} competition track with the Puzzlebot
           AMR. Trucks are at the far end; racks (two levels) occupy the
           centre; roller tables are at the near end.}
  \label{fig:pista}
\end{figure}

\subsection{Objectives and Contributions}

The project pursued one functional goal --- complete a full
\textit{pick-and-place} mission autonomously --- decomposed into nine
engineering objectives, one per subsystem (localization, navigation,
mission orchestration, lifting, perception, voice, dashboard, digital
twin, and safety). Relative to a baseline that would have stitched
together off-the-shelf ROS~2 packages, the principal contributions of
this work are:

\begin{enumerate}[noitemsep, leftmargin=*]
  \item A \textbf{from-scratch C++ Monte-Carlo localization} node with an
        optional CUDA likelihood kernel, ICP refinement, and a
        graph-SLAM back-end for loop closure --- a drop-in replacement
        for \texttt{slam\_toolbox}~+~\texttt{amcl}.
  \item An \textbf{ArUco beacon localization layer} that acts as a
        rescue seed for the particle filter, eliminating the orientation
        drift that featureless corridors induce in pure laser MCL.
  \item A \textbf{self-contained navigation stack} (A*, Bug1,
        pure-pursuit, local costmap) that replaces Nav2 while remaining
        small enough to read and tune by hand.
  \item A \textbf{deterministic mission orchestrator} built on YASMIN
        with a live JSON debugger for field bring-up.
  \item A \textbf{custom voice front-end} (MFCC~+~VQ~+~HMM) that runs on
        CPU with no pre-trained network and no cloud dependency.
\end{enumerate}

\subsection{Report Organization}

Section~\ref{sec:arch} describes the distributed deployment and the
ROS~2 graph. Section~\ref{sec:subsys} details each subsystem --- SLAM,
navigation, mission control, lifting, perception, voice, dashboard,
digital twin, and the industrial-standards mapping. Section~\ref{sec:challenges}
collects the integration problems that were diagnosed and fixed,
Section~\ref{sec:results} reports the end-to-end mission results, and
Section~\ref{sec:status} gives the final task backlog.

% ════════════════════════════════════════════════════════════════════════════
\section{Final System Architecture}
\label{sec:arch}

\subsection{Distributed Deployment}

The system runs as two launch groups on the same \texttt{ROS\_DOMAIN\_ID}:

\begin{itemize}[noitemsep, leftmargin=*]
  \item \textbf{Jetson Nano} (\texttt{robot.launch.py}): LiDAR driver,
        \texttt{real\_odom} (differential kinematics + TF),
        \texttt{slam\_node} (C++ MCL), \texttt{aruco\_localization},
        \texttt{lifting\_node} (SPI → FPGA), \texttt{qr\_quad\_alignment},
        \texttt{logo\_stop\_debug}, \texttt{twist\_relay}.
  \item \textbf{Laptop} (\texttt{laptop.launch.py}): \texttt{nav\_node}
        (A* + Bug1 + pure-pursuit), \texttt{state\_machine\_node}
        (YASMIN), \texttt{dashboard\_node} (Flask), \texttt{voice\_node}
        (HMM), RViz~2, Gazebo digital twin.
\end{itemize}

\paragraph{Anti-brownout relay} all upstream nodes publish velocity
commands to \texttt{/cmd\_vel\_in}; the \texttt{twist\_relay} node
applies a ramp filter and forwards to \texttt{/cmd\_vel} → micro-ROS →
motors. This prevents voltage-induced Jetson resets from sudden current
spikes.

\begin{table}[h]
\caption{ROS~2 packages in the final system}
\label{tab:packages}
\renewcommand{\arraystretch}{1.2}
\centering
\begin{tabular}{>{\raggedright}p{2.2cm}
                >{\centering}p{1.0cm}
                >{\raggedright\arraybackslash}p{4.2cm}}
\toprule
\textbf{Package} & \textbf{Lang} & \textbf{Purpose} \\
\midrule
\texttt{slam}          & C++    & MCL/AMCL + CUDA + graph SLAM \\
\texttt{localization}  & C++    & EKF + ICP velocity bridge \\
\texttt{controller}    & Python & Kinematics, odometry, twist relay \\
\texttt{navigation}    & Python & A* + Bug1 + pure-pursuit \\
\texttt{mission\_control} & Python & YASMIN state machine \\
\texttt{perception}    & Python & ArUco, QR dock, logo YOLO \\
\texttt{lifting}       & Python & HAL SPI → FPGA \\
\texttt{voice\_control} & Python & MFCC + VQ + HMM \\
\texttt{dashboard}     & Python & Flask + SocketIO \\
\texttt{description}   & CMake  & URDF, Gazebo world, digital twin \\
\texttt{bringup}       & CMake  & Launch files \\
\bottomrule
\end{tabular}
\end{table}

\subsection{Coordinate Frames and the TF Tree}

All geometry is expressed in a standard REP-105 frame chain. The
\texttt{map} frame is the static world anchored to the saved occupancy
grid; \texttt{odom} is a smooth, drift-prone frame produced by wheel
odometry; \texttt{base\_link} is the robot body; and \texttt{laser} and
\texttt{camera\_link} are the sensor frames fixed by the URDF. The
localization stack publishes the \texttt{map}~$\rightarrow$~\texttt{odom}
correction, while \texttt{real\_odom} publishes
\texttt{odom}~$\rightarrow$~\texttt{base\_link}:

\[
T_{\text{map}}^{\text{base}} =
\underbrace{T_{\text{map}}^{\text{odom}}}_{\text{SLAM/MCL}}
\cdot
\underbrace{T_{\text{odom}}^{\text{base}}}_{\text{wheel odometry}}.
\]

Decoupling the two transforms this way is what lets MCL apply a discrete
pose correction (which would otherwise teleport the robot) without ever
introducing a discontinuity into the high-rate \texttt{odom} frame that
the controller integrates.

\subsection{Data Flow}

Sensor topics (\texttt{/scan}, \texttt{/cam\_img}, the encoder
velocities) originate on the Jetson and are consumed locally by SLAM and
perception; only compact estimates (\texttt{/slam\_pose}, \texttt{/map},
\texttt{/aruco\_pose\_estimate}, \texttt{/logo\_detection}) cross the
Wi-Fi link to the laptop, keeping the wireless DDS traffic well below the
bandwidth that raw images would demand. Command topics flow the other
way: the planner and the state machine publish to \texttt{/cmd\_vel\_in},
the relay ramps them, and \texttt{/cmd\_vel} is bridged to the motors
through micro-ROS. Figure~\ref{fig:arch} shows the full graph.

\begin{figure}[htbp]
  \centering
  \includegraphics[width=\columnwidth]{assets/architecture.png}
  \caption{End-to-end ROS~2 data flow. On-board nodes (Jetson) handle
           sensing, SLAM, perception, and lifting; off-board nodes
           (laptop) handle planning, mission control, the dashboard, and
           the digital twin. Only compact estimates cross the Wi-Fi/DDS
           boundary.}
  \label{fig:arch}
\end{figure}

% ════════════════════════════════════════════════════════════════════════════
\section{Subsystem Descriptions}
\label{sec:subsys}

\subsection{M3. SLAM and Localization}

\paragraph{C++ MCL/AMCL node}
\texttt{slam\_node} (C++) runs the full particle-filter loop at 10~Hz
on every \texttt{/scan} arrival:

\begin{enumerate}[noitemsep, leftmargin=*]
  \item \textbf{Predict} --- wheel odometry + ICP scan-to-scan as a
        slip-robust rotation prior.
  \item \textbf{Update} --- likelihood-field sensor model scores each
        particle against the wall map; a CUDA kernel accelerates this
        step when the GPU is available.
  \item \textbf{Resample} --- low-variance resampling with recovery
        (random injection) triggered when \texttt{neff} falls below
        15\% of 600 particles.
  \item \textbf{Scan-to-map refine} --- ICP over a 4~m search radius,
        0.3~m maximum correction, with a 15° coarse yaw sweep.
\end{enumerate}

A back-end thread runs \textbf{graph SLAM}: it accumulates keyframes
and detects loop closures, then re-rasterizes the map. Two operating
modes are provided: \texttt{mapping} (build a new \texttt{.pgm}) and
\texttt{navigation} (load the saved map, MCL only).

Map parameters: 0.04~m cells, 400 × 400 grid (16 × 16~m), incremental
log-odds (occupied $+0.5$, free $-0.2$).

\paragraph{Motion model}
Each particle $i$ is propagated with the odometry sample model: the
wheel displacement $(\delta_{\text{trans}}, \delta_{\text{rot1}},
\delta_{\text{rot2}})$ is perturbed by zero-mean Gaussian noise whose
variance scales with the motion magnitude,
\[
\hat\delta_{\text{rot}} = \delta_{\text{rot}}
  - \mathcal{N}\!\left(0,\; \alpha_1 \delta_{\text{rot}}^2
                            + \alpha_2 \delta_{\text{trans}}^2\right),
\]
so that fast or sharply-turning motion inflates the cloud and slow
straight motion keeps it tight. The ICP scan-to-scan estimate replaces
the raw wheel rotation whenever the two disagree by more than a
threshold, which makes the prior robust to wheel slip on the track's
smooth surface.

\paragraph{Sensor model}
The likelihood field pre-computes, for every free cell, the Euclidean
distance $d$ to the nearest occupied cell. A beam endpoint falling on a
cell of distance $d$ contributes
\[
p(z\mid x) \;\propto\;
  z_{\text{hit}}\, e^{-d^2 / 2\sigma^2} \;+\; z_{\text{rand}},
\]
and the per-particle weight is the product over a subsampled set of beams
(every $k$-th ray) for speed. The CUDA kernel evaluates all 600
particles against all beams in parallel; on the Jetson's GPU this cut the
update step from the dominant cost of the loop to a negligible fraction,
which is what allowed the full filter to hold 10~Hz.

\paragraph{ArUco rescue beacons}
\texttt{aruco\_localization} (Python, \texttt{perception} package)
operates 20 fixed floor markers (dictionary \texttt{DICT\_ARUCO\_ORIGINAL},
9~cm side). For each detection it computes the robot's absolute pose
$T_{\text{map}}^{\text{base}} = T_{\text{map}}^{\text{marker}}
\cdot (T_{\text{cam}}^{\text{marker}})^{-1}
\cdot (T_{\text{base}}^{\text{cam}})^{-1}$
and publishes to \texttt{/aruco\_pose\_estimate}
(\texttt{PoseWithCovarianceStamped}).
When two or more markers are visible simultaneously, a
\textbf{2-D Umeyama rigid registration} fits the observed positions to
the canonical map, rejecting outliers greedily; baseline
$\geq$0.25~m and fit RMSE $\leq$0.10~m are required to publish.
\texttt{slam\_node} consumes the estimate as a rescue seed: if the MCL
belief disagrees by more than \texttt{aruco\_snap\_tol}, the particle
cloud is re-centred there and the laser snaps it to the walls.

A visual web editor (\texttt{perception/tools/aruco\_map\_editor.py})
allows the map to be updated interactively at
\texttt{http://127.0.0.1:8770}.

\begin{figure}[htbp]
  \centering
  \includegraphics[width=\columnwidth]{assets/mapa_montecarlo.jpeg}
  \caption{RViz~2 view of the final localization stack: occupancy grid,
           MCL particle cloud (arrows), RPLidar A1 scan overlay, and
           the ArUco detection debug image (inset top-right) showing
           detected marker IDs and their estimated poses.}
  \label{fig:slam}
\end{figure}

\begin{figure}[htbp]
  \centering
  \includegraphics[width=\columnwidth]{assets/aruco_gazebo_rviz.jpeg}
  \caption{ArUco beacon localization, side by side. Left: the robot in
           the Gazebo digital twin observing the floor markers. Right: the
           corresponding RViz~2 view with the detected marker poses and
           the resulting rescue seed used to re-centre the MCL particle
           cloud.}
  \label{fig:aruco}
\end{figure}

\subsection{MR. Navigation}

\texttt{nav\_node} loads the pre-saved map (\texttt{.pgm} +
\texttt{.yaml}) and \texttt{waypoints.yaml} at startup.
Named zones follow the convention \texttt{truck\_*},
\texttt{rack\_*}, \texttt{roller\_*}.

\paragraph{Global planner --- A*}
Plans over an inflated costmap (default 0.30~m, hard floor 0.15~m) on an
8-connected grid. Each node expands with cost
\[
f(n) = g(n) + h(n), \qquad
h(n) = \lVert n - \text{goal} \rVert_2,
\]
where $g(n)$ accumulates diagonal-aware step costs plus the inflation
penalty of each traversed cell, and the Euclidean heuristic $h(n)$ keeps
the search admissible (it never over-estimates) so the returned path is
optimal for the chosen costmap. An inflation \textbf{fallback ladder}
(0\%, 25\%, 50\%, 75\%) retries with reduced margins if the goal cell is
occupied after inflation, so the robot always attempts a plan even when a
pallet has been parked tight against a wall.

\paragraph{Path follower --- pure-pursuit PD}
Tracks the planned path at $v = 0.10$~m/s with a 0.50~m look-ahead. Given
the look-ahead point at lateral offset $y$ and range $L_d$ in the robot
frame, the pursuit curvature is
\[
\kappa = \frac{2y}{L_d^{2}}, \qquad
\omega = \operatorname{clip}\!\big(K_p\,\kappa\,v + K_d\,\dot e,\;
                                   \pm\,\omega_{\max}\big),
\]
with a low-pass derivative filter ($\tau = 0.2$~s) on the heading error
$e$ to damp oscillation. A \textbf{reverse maneuver} is triggered when the
look-ahead point lies behind the robot and within 0.6~m, allowing
recovery from over-shoots instead of a wide loop-around.

\begin{figure}[htbp]
  \centering
  \includegraphics[width=\columnwidth]{assets/nav_path.png}
  \caption{Navigation in action: the A* global plan, the pure-pursuit
           look-ahead target, and a Bug1 boundary-following detour when a
           previously unmapped obstacle blocks the planned path.}
  \label{fig:nav}
\end{figure}

\paragraph{Reactive avoidance --- Bug1 + local costmap}
On obstacle hit the robot switches to Bug1 wall-follow (target distance
0.45~m, right-hand rule). The leave-point is tracked; three watchdogs
bound the maneuver: 25~s in-vicinity, 90~s absolute, and 15~s
no-progress. The \textbf{local costmap} (\texttt{local\_costmap.py})
provides a 2~m rolling patch with exponential decay ($\tau = 0.7$~s),
giving reactive avoidance a temporally-smoothed obstacle signal.

Key tuned parameters: $K_p = 1.5$, $K_d = 0.5$,
$\omega_{\max} = 0.21$~rad/s; goal tolerance 0.03~m,
heading tolerance 0.04~rad; approach speed floor 0.04~m/s.

\subsection{MR. Mission Control --- YASMIN State Machine}

The \texttt{state\_machine\_node} hosts a YASMIN
finite-state machine with eight states:

\begin{align*}
\text{IDLE}
  &\xrightarrow{\text{mission\_received}} \text{SEARCH} \\
  &\xrightarrow{\text{aligned}} \text{PICK / PICK\_FROM\_RACK} \\
  &\xrightarrow{\text{picked}} \text{NAV\_TO\_TRUCK} \\
  &\xrightarrow{\text{arrived}} \text{RELEASE\_LOAD} \\
  &\xrightarrow{\text{released}} \text{MISSION\_DONE} \to \text{IDLE}
\end{align*}

Any state can transition to \texttt{MISSION\_FAILED} on timeout or
stop command, which resets back to \texttt{IDLE}.

\paragraph{Mission types} two mission types are supported:
\begin{itemize}[noitemsep, leftmargin=*]
  \item \textbf{Rollers → Truck (Mission~1):} the robot navigates to
        the roller staging area, aligns using the E80 logo classifier
        for a vision stop, picks the pallet, then navigates to the
        correct truck (identified by logo: amazon/pepsi/walmart) and
        releases.
  \item \textbf{Racks → Truck (Mission~2):} same flow but the pick is
        from a two-level rack (\texttt{PICK\_FROM\_RACK} state); the
        lifter reaches level~5 for the upper shelf.
\end{itemize}

\paragraph{Shared blackboard}
States do not call each other directly; they communicate through a single
YASMIN blackboard whose keys are written by ROS~2 callbacks and read by
the state \texttt{execute} methods (Table~\ref{tab:blackboard}). This
keeps the FSM logic pure --- a state's outcome is a function of the
blackboard snapshot alone --- which is what makes the \texttt{step} and
\texttt{force\_outcome} debug commands deterministic and safe to use mid
mission.

\begin{table}[h]
\caption{Principal YASMIN blackboard keys}
\label{tab:blackboard}
\renewcommand{\arraystretch}{1.2}
\centering
\begin{tabular}{>{\raggedright}p{2.3cm}
                >{\raggedright\arraybackslash}p{5.1cm}}
\toprule
\textbf{Key} & \textbf{Meaning / writer} \\
\midrule
\texttt{mission}      & active mission JSON (source, dest, type) \\
\texttt{robot\_pose}  & latest \texttt{/slam\_pose} sample \\
\texttt{aligned}      & logo/QR alignment converged flag \\
\texttt{lifter\_done} & last \texttt{/lifter\_level} ack \\
\texttt{stop\_flag}   & protective-stop latch (voice/dashboard) \\
\texttt{step\_mode}   & debugger single-step gate \\
\bottomrule
\end{tabular}
\end{table}

\paragraph{Live debugger}
A \texttt{/sm/control} topic accepts JSON commands at runtime:
\texttt{pause}, \texttt{resume}, \texttt{step},
\texttt{set\_step\_mode}, \texttt{force\_outcome}, \texttt{abort}.
Every transition is published to \texttt{/sm/transition}; the full
blackboard is snapshotted periodically to \texttt{/sm/blackboard}, which
the dashboard renders live so an operator can watch the mission advance
state by state.

\subsection{M2. FPGA Lifter --- SPI HAL}

The \texttt{lifting\_node} translates \texttt{/lifter\_level} (UInt8,
range 0–5) into SPI commands to the Tang~Nano~20K FPGA
(bus~0, device~0, 1~MHz, mode~0). The FPGA maps each 3-bit binary
level to a servo angle:

\begin{table}[h]
\caption{Lifter level encoding (SPI → FPGA → servo)}
\label{tab:lifter_encoding}
\renewcommand{\arraystretch}{1.2}
\centering
\begin{tabular}{cccc}
\toprule
\textbf{Level} & \textbf{Binary} & \textbf{Servo} & \textbf{Use} \\
\midrule
0 & 000 & 95° & Rest / floor transport \\
1 & 001 & 70° & Low travel height \\
2 & 010 & 65° & Intermediate \\
3 & 011 & 40° & Pick pallet on floor \\
4 & 100 & 25° & Carry / transport \\
5 & 101 &  0° & Pick pallet rack level~2 \\
\bottomrule
\end{tabular}
\end{table}

A pluggable HAL abstraction (\texttt{lifting/hal/}) provides
\texttt{JetsonHAL} (SPI via \texttt{spidev}), \texttt{MockHAL}
(simulation), and the legacy \texttt{JetsonGPIOHAL} (3 GPIO pins).
The active HAL is selected by \texttt{lifting\_params.yaml}.
Level state is cached to \texttt{/tmp} so a mid-mission Jetson restart
does not lose the current fork height.

\begin{figure}[htbp]
  \centering
  \fbox{\parbox{0.9\columnwidth}{\centering\vspace{6pt}
    \textbf{Final demonstration video} \\[6pt]
    Full pick-and-place run: Mission~1 (Rollers → Truck) and
    Mission~2 (Racks → Truck) on the physical robot.
    Shows SLAM, ArUco rescue, navigation, logo stop, pick, and
    delivery. \\[8pt]
    \href{https://drive.google.com/drive/folders/1t7DX9yP-UUfOrtskFwrFIzdNqKrJuPC_?usp=sharing}{%
      \texttt{Two Missions}} \\[6pt]
  }}
  \caption{Final demonstration: complete autonomous missions on
           hardware.}
  \label{fig:demo_video}
\end{figure}

\subsection{M2. Hardware and Electronics}

The robot is a differential-drive Puzzlebot. Two DC gear-motors are
driven through an H-bridge whose direction and PWM lines are commanded by
a micro-ROS firmware that subscribes to \texttt{/cmd\_vel} and closes a
low-level velocity loop on the wheel encoders (\texttt{/VelocityEncR},
\texttt{/VelocityEncL}). The differential kinematics that the controller
inverts to produce wheel set-points are
\[
v = \frac{r}{2}(\omega_R + \omega_L), \qquad
\omega = \frac{r}{b}(\omega_R - \omega_L),
\]
for wheel radius $r$ and track width $b$. Sensing is an RPLidar~A1
(\texttt{/scan}) and a forward camera (\texttt{/cam\_img}); the lifter is
a servo driven by the Tang~Nano~20K FPGA.

A recurring failure mode early on was a \textbf{voltage brownout}: a
sudden velocity step drew a current spike through the H-bridge that
dropped the rail below the Jetson's reset threshold, rebooting the
on-board stack mid mission. The \texttt{twist\_relay} ramp filter
(Section~\ref{sec:arch}) bounds $\lvert\Delta v\rvert$ per control tick
and was the single change that made long autonomous runs reliable.

\begin{figure}[htbp]
  \centering
  \includegraphics[width=\columnwidth]{assets/hardware.jpeg}
  \caption{The Puzzlebot AMR platform: Jetson Nano compute, RPLidar~A1,
           forward camera, H-bridge motor driver, and the Tang~Nano~20K
           FPGA servo lifter.}
  \label{fig:hardware}
\end{figure}

\subsection{MR. Visual Alignment and Perception}

\paragraph{E80 logo classifier (\texttt{logo\_classifier.py})}
A YOLOv8 model (\texttt{weights.pt}) runs at 15~Hz, detects the three
company logos (amazon, pepsi, walmart) on the trucks and pallets, sorts
detections left-to-right by bounding-box centre, and publishes the
ordered list as JSON to \texttt{/logo\_detection}. Confidence threshold
0.35.

During \texttt{SEARCH} the state machine uses the logo stop to halt
precisely in front of the pallet; during \texttt{RELEASE\_LOAD} it
selects the correct truck by matching the pallet's QR-decoded company
name (with spelling-insensitive aliases: WOLMAR~→~walmart,
EMEZON~→~amazon, POPSI~→~pepsi) to the logo observed on each dock.

\paragraph{QR docking (\texttt{qr\_quad\_alignment.py})}
The parabola + wheel-odometry strategy:
(1)~\textbf{PLAN} --- at detection, snapshot the QR pose and plan a
parabola $y = ax^2$ in the robot frame; reset wheel odometry;
(2)~\textbf{FOLLOW} --- pure pursuit on the parabola using only
\texttt{odometry\_tracker.py}, immune to occlusion while the robot's
body blocks the marker;
(3)~\textbf{DOCK} --- final straight approach with QR feedback if
visible, open-loop otherwise.

\begin{figure}[htbp]
  \centering
  \includegraphics[width=\columnwidth]{assets/qr_docking.jpeg}
  \caption{QR docking. The robot snapshots the QR pose, plans a
           parabola $y = ax^2$ in its own frame, and follows it on wheel
           odometry --- staying aligned even when its chassis occludes the
           marker during the final approach.}
  \label{fig:qr_docking}
\end{figure}

\paragraph{ArUco detection (\texttt{aruco\_detector.py})}
Compatible with OpenCV~4.6 and 4.7+ via a version-detection shim.
Publishes 3-D poses (\texttt{rvec} + \texttt{tvec} → quaternion) to
\texttt{/aruco\_poses} (\texttt{PoseArray}) and IDs to
\texttt{/aruco\_ids}.

\begin{figure}[htbp]
  \centering
      \includegraphics[width=\columnwidth]{assets/perception.jpeg}
  \caption{Perception during \texttt{RELEASE\_LOAD}: the YOLOv8 classifier
           localizes the three dock logos while the QR decoder reads the
           pallet's company tag for truck selection.}
  \label{fig:perception}
\end{figure}

\subsection{M4. Voice Control}

The \texttt{voice\_node} implements a fully custom speech-recognition
pipeline with no pre-trained models:

\begin{align*}
\text{Audio}_{16\,\text{kHz}}
  &\xrightarrow{\text{VAD}} \xrightarrow{\text{MFCC}} \\
  &\xrightarrow{\text{VQ}_{K=256}} \xrightarrow{\text{HMM}}
  \text{command}
\end{align*}

VAD uses an adaptive RMS threshold (0.015); MFCC features are
quantised against per-word codebooks (VQ, $K = 256$ centres) and
then scored by per-word left-to-right HMMs. The word with the highest
log-likelihood wins.

\paragraph{Feature extraction}
Each VAD-gated utterance is framed (25~ms windows, 10~ms hop) and reduced
to MFCCs: the magnitude spectrum is mapped through a mel filterbank and
the log-energies are decorrelated by a DCT,
\[
c_n = \sum_{m=1}^{M} \log E_m \,
        \cos\!\Big[\tfrac{\pi n}{M}\big(m - \tfrac12\big)\Big],
\]
keeping the low-order coefficients $c_n$ as the per-frame feature.

\paragraph{Recognition}
Each frame vector is replaced by its nearest VQ codeword, turning the
utterance into a discrete symbol sequence $O = o_1\dots o_T$. A separate
left-to-right HMM $\lambda_w$ is trained per vocabulary word, and the
recognized word maximizes the forward likelihood
\[
\hat w = \arg\max_w \; P(O \mid \lambda_w),
\]
evaluated with the forward algorithm. Because every word has its own
small model and the front-end is plain DSP, the whole recognizer runs in
real time on the laptop CPU with no GPU and no pre-trained network ---
the explicit design goal of avoiding Whisper.

Commands are parsed into three categories:
\begin{itemize}[noitemsep, leftmargin=*]
  \item \textbf{Mission} (\texttt{start, stop, break, resume,
        continue, leave}) --- published to \texttt{/mission} or the
        YASMIN blackboard.
  \item \textbf{Control} (\texttt{lift, lower}) --- published to
        \texttt{/lifter\_level}.
  \item \textbf{Motion} (\texttt{forward, backward, left, right}) ---
        published at 0.1~Hz to \texttt{/cmd\_vel\_in} during the
        active window.
\end{itemize}

\subsection{MR. Web Dashboard}

The Flask + SocketIO dashboard (port 8080) provides:

\begin{itemize}[noitemsep, leftmargin=*]
  \item \textbf{Live camera stream} (MJPEG) from \texttt{/cam\_img}.
  \item \textbf{SLAM map canvas} --- \texttt{ros\_bridge.py} encodes
        each \texttt{/map} (\texttt{OccupancyGrid}) update as PNG and
        pushes it via SocketIO; the robot pose is overlaid.
  \item \textbf{Telemetry panel} --- pose (\texttt{/slam\_pose}),
        YASMIN state (\texttt{/robot\_state}), active mission, lifter
        level.
  \item \textbf{Mission control} --- dispatch missions (type, source,
        destination), abort / pause / resume the state machine.
  \item \textbf{Waypoint capture} and \textbf{map save} buttons.
  \item \textbf{Voice-command panel} --- browser microphone → WebM →
        \texttt{pydub} decode → HMM recognizer → \texttt{/voice\_command}.
  \item \textbf{Safety watchdog} --- 350~ms dead-man's switch publishes
        zero velocity if no teleoperation arrives.
\end{itemize}

\begin{figure}[htbp]
  \centering
  \includegraphics[width=\columnwidth]{assets/dashboard.jpeg}
  \caption{Flask + SocketIO web dashboard (port 8080).
           Panels (left to right): SLAM occupancy-grid map with
           robot pose overlay; MJPEG camera feed; telemetry (YASMIN
           state, pose, lifter level); mission dispatch and control
           buttons.}
  \label{fig:dashboard}
\end{figure}

\subsection{MR. Digital Twin}

The Gazebo digital twin is not a physics simulation; it is a live
\textbf{environment visualizer} that mirrors the physical robot's
estimated pose in real time:

\begin{itemize}[noitemsep, leftmargin=*]
  \item \texttt{gemelo\_mirror} subscribes to \texttt{/slam\_pose} and
        teleports the virtual robot via Ignition's \texttt{set\_pose}
        service at every SLAM update.
  \item The \texttt{almacen\_racks.world} scene contains measured-to-reality
        rack and roller positions (\texttt{config/rack\_groups.yaml},
        \texttt{roller\_groups.yaml}); the layout is regenerated from
        \texttt{scripts/gen\_warehouse.py}.
  \item Textured ArUco markers match the physical placement so the
        visualizer can be overlaid on \texttt{aruco\_map.yaml}.
\end{itemize}

\begin{figure}[htbp]
  \centering
  \includegraphics[width=\columnwidth]{assets/gemelo.jpeg}
  \caption{Live Gazebo digital twin. The virtual Puzzlebot tracks the
           real robot's \texttt{/slam\_pose} in real time via the
           \texttt{gemelo\_mirror} node. Racks, rollers, walls, and
           textured ArUco markers match the physical track layout.}
  \label{fig:gemelo}
\end{figure}

\subsection{M1. Industrial Standards}

\paragraph{ISO~10218-1/2 (Industrial robot safety)}
\begin{itemize}[noitemsep, leftmargin=*]
  \item \textbf{Speed and separation:} maximum translational velocity
        capped at 0.10~m/s (nav\_params); Bug1 maintains a 0.45~m
        wall-separation distance.
  \item \textbf{Protective stop:} \texttt{stop} voice token and
        dashboard abort button publish zero velocity and raise
        \texttt{stop\_flag} in the YASMIN blackboard within one control
        cycle ($< 100$~ms).
  \item \textbf{Anti-brownout relay:} the \texttt{twist\_relay} ramp
        prevents abrupt current spikes that could damage the H-bridge.
\end{itemize}

\paragraph{ISO~13482 (Personal care robots --- safety)}
The OWL ontology (Sprint~3) was extended with
\texttt{HazardZone} and \texttt{SafetyConstraint} classes linked to
the robot's workspace zones. The dead-man's switch in the dashboard
(350~ms watchdog) fulfils the unintended-motion requirement for a
teleoperated service robot.

% ════════════════════════════════════════════════════════════════════════════
\section{Challenges Overcome}
\label{sec:challenges}

\begin{itemize}[leftmargin=*]
  \item \textbf{Real-world SLAM smearing:} three independent root causes
        (DDS multicast collision, clock offset rejecting TF frames, QoS
        mismatch silently dropping scans) all produced the same symptom.
        Each was diagnosed and fixed: \texttt{IGN\_IP=localhost}, wall-clock
        time with 0.3~s transform tolerance, and matched QoS profiles.
  \item \textbf{MCL drift in low-feature corridors:} the aisle between
        racks has parallel featureless walls that make the particle filter
        lose orientation. The ArUco beacon layer resolves this: a single
        visible marker re-seeds the belief, and the laser immediately snaps
        the particles to the walls.
  \item \textbf{PnP ambiguity for floor markers:} standard OpenCV
        \texttt{solvePnP} returns two flip solutions for upward-facing
        markers. \texttt{aruco\_localization} tests four in-plane
        conventions and selects the solution with maximum Z-axis
        uprightness, making the estimate immune to the floor-marker
        ambiguity.
  \item \textbf{Logo mis-spelling in QR payloads:} physical QR codes
        encode company names with inconsistent spelling (``Wolmar'',
        ``Emezon'', ``Popsi''). The \texttt{logo\_aliases} map in
        \texttt{zones.yaml} normalises all variants case- and
        punctuation-insensitively before truck selection.
  \item \textbf{Vision occlusion during docking:} the robot's chassis
        blocks the QR marker at close range. The parabola + odometry
        strategy (\texttt{qr\_quad\_alignment.py}) eliminates the
        dependency on continuous vision during the FOLLOW phase.
  \item \textbf{Voltage brownout on acceleration:} rapid velocity
        commands caused H-bridge current spikes that reset the Jetson.
        The \texttt{twist\_relay} ramp filter limits $\Delta v$ per tick.
  \item \textbf{YASMIN + ROS~2 thread safety:} the state machine runs in
        a background thread while ROS~2 callbacks update the blackboard
        from the executor thread. CPython's GIL guarantees atomicity for
        the scalar and dict types used on the blackboard for this use case.
\end{itemize}

% ════════════════════════════════════════════════════════════════════════════
\section{Results}
\label{sec:results}

Both target missions were executed successfully end-to-end on the
physical robot:

\begin{itemize}[noitemsep, leftmargin=*]
  \item \textbf{Mission~1 --- Rollers → Truck:} robot navigates to the
        roller staging area, uses the E80 logo classifier for a vision
        stop, picks the pallet at floor level (lifter level~3), navigates
        to the correct truck identified by logo matching, and deposits the
        pallet.
  \item \textbf{Mission~2 --- Racks → Truck:} same pipeline with the
        \texttt{PICK\_FROM\_RACK} state; the lifter reaches level~5 (0°)
        to engage pallets on the upper shelf.
\end{itemize}

Table~\ref{tab:metrics} summarizes the operating characteristics observed
during the validation sessions. The headline result is qualitative ---
\textbf{both missions complete autonomously, repeatably} --- but the
quantitative figures below capture the tuning targets the stack was
brought to. Figure~\ref{fig:mission_seq} shows a representative run.

\begin{table}[h]
\caption{Observed system characteristics}
\label{tab:metrics}
\renewcommand{\arraystretch}{1.2}
\centering
\begin{tabular}{>{\raggedright}p{4.4cm}
                >{\raggedleft\arraybackslash}p{2.8cm}}
\toprule
\textbf{Quantity} & \textbf{Value} \\
\midrule
MCL / SLAM update rate            & 10~Hz \\
Particle count                    & 600 \\
Logo classifier rate              & 15~Hz \\
Nominal cruise speed              & 0.10~m/s \\
Goal position tolerance           & 0.03~m \\
Goal heading tolerance            & 0.04~rad \\
Bug1 wall-separation distance     & 0.45~m \\
Protective-stop latency           & $<100$~ms \\
Missions validated on hardware    & 2 / 2 \\
\bottomrule
\end{tabular}
\end{table}

\begin{figure}[htbp]
  \centering
  \includegraphics[width=\columnwidth]{assets/mission_seq.png}
  \caption{Representative Mission~1 sequence on hardware: navigation to
           the rollers, logo-based vision stop, pallet pick, transport,
           truck alignment, and release.}
  \label{fig:mission_seq}
\end{figure}

A demonstration video of both missions is linked in
Figure~\ref{fig:demo_video}.

% ════════════════════════════════════════════════════════════════════════════
\section{Conclusion}

The Puzzlebot AMR completed both target missions autonomously on the
physical hardware. The fully custom stack --- C++ MCL SLAM with CUDA
scoring and graph-SLAM loop closure, ArUco beacon localization as an
MCL rescue layer, A*~+~Bug1~+~pure-pursuit navigation, a YASMIN
orchestrator, YOLOv8 logo-based truck selection, QR docking via
parabola~+~odometry, an MFCC~+~VQ~+~HMM voice pipeline, a live digital
twin, and a Tang~Nano~20K SPI lifter --- demonstrates that the entire
AMR autonomy stack can be implemented from first principles without
relying on Nav2, Whisper, or \texttt{robot\_localization}. All 23
planned tasks were completed before the final demonstration. The modular
package structure and distributed Jetson~+~laptop deployment proved
robust throughout integration testing and the final demonstration.

\end{document}
